\scanpage{6}
À vrai dire, l'\deleted{application}\inserted{injectivité} de l'application sur
les sommets, alors que l'application \deleted{\unsure{sur}} des
côtés comb vers les côtés géom peut
n'être pas injective, est \unsure{anachronique|encore que}
dans le cas non propre [\deleted{\unsure{en|on|où}} \deleted{dans le cas propre} la donnée
du polygône géom \inserted{propre} équivaut à la donnée
de l'ens des sommets et de l'ens
des côtés, satisfaisant les conditions suivantes\,: 
tt sommet incident à ex. deux côtés, tt
\unsure{côtés|côté} incident à exactement une \unsure{arête},
\deleted{\unsure{tout}} plus condition de connexité combinatoire\ldots]
La bonne notion est
donc celle de \uline{repr.} \inserted{(ou \uline{réalisation})} \uline{géom} d'un polyg.
combinatoire, qui revient à la donnée
de \uline{deux} applications -- sur l'ens des
sommets \uline{et} sur l'ens des droites --
respectant les relations d'incidence\ldots.
(Lorsque deux sommets incidents ont
\deleted{\unsure{l'image des pt de vue sommets}}
\deleted{\unsure{l'un des sommets a}} \inserted{\unsure{tjrs}} des images
\deleted{\unsure{pour être un polygône|pour être un polygone}} distinctes
par l'application $\varphi_{\mathrm{sommets}}$, cette dernière
suffit à elle seul à déterminer la réalisation,
et elle peut être choisie arbitrairement,
% archive sheet number: 5

